\section{Exact cofactor descent on a global prime band}
\label{sec:tpc45-cofactor-descent}

We begin with a prime band defined from the whole equality packet, not
from one selected row.  This distinction is necessary because the sign
attached to a prime is global: the same prime cannot be frozen in one
occurrence and left random in another.
In this section and in the determinant-ladder discussion we denote a
band prime by \(q\); the generic incidence and Gram sections denote the
same prime role by \(p\).

Let \(\mathfrak A_X\) be the set of nonzero equality atoms.  An atom
has a coordinate \(i\), whose data include an orbit time \(j=j(i)\),
and a source row
\(\alpha=(\ell_\alpha,d_\alpha)\).  Put
\begin{equation}
 n_{\alpha,j}=\ell_\alpha d_\alpha j+h,
 \qquad
 u_{\alpha,j}=d_\alpha n_{\alpha,j}.
 \label{eq:tpc45-n-u}
\end{equation}
On the nonzero support, \(d_\alpha\) and \(n_{\alpha,j}\) are
squarefree and
\begin{equation}
 (d_\alpha,n_{\alpha,j})=(d_\alpha,h)=1.
 \label{eq:tpc45-source-target-coprime}
\end{equation}
Use the global target endpoints \(N_\pm\) from the actual incidence
interface, chosen as the extrema of the nonzero support, and define
\begin{equation}
 \ell_+:=\max_\alpha\ell_\alpha.
 \label{eq:tpc45-support-extrema}
\end{equation}
Empty packets are harmless, so we suppose \(\mathfrak A_X\ne\varnothing\).
The fixed dyadic supports give
\begin{equation}
 N_-=X^{1+o(1)},
 \qquad N_+=X^{1+o(1)},
 \qquad
 \ell_+=X^{\lambda_L+o(1)},
 \qquad
 \lambda_L=\frac{99979}{210000}.
 \label{eq:tpc45-extremal-scales}
\end{equation}

\begin{definition}[Global cofactor band]
\label{def:tpc45-global-cofactor-band}
Let \(R=R_X\) satisfy
\begin{equation}
 1\le R<\frac{N_-}{\sqrt{N_+}},
 \label{eq:tpc45-R-range}
\end{equation}
and set
\begin{equation}
 \cP_R:=
 \left\{q\in\mathbb P:
  \sqrt{N_+}<q\le \frac{N_-}{R}
 \right\}.
 \label{eq:tpc45-global-cofactor-band}
\end{equation}
Every sign \(\varepsilon_q\), \(q\in\cP_R\), will be set to
the physical value \(-1\) in every occurrence of \(q\).
\end{definition}

If \(R=X^{\theta+o(1)}\), \(0\le\theta<1/2\), the band is
asymptotically nonempty and has endpoint exponents \(1/2\) and
\(1-\theta\).  The clean balanced choice \(R=\ell_+\) will be denoted by
\begin{equation}
 \cP_{\rm bal}:=\cP_{\ell_+}.
 \label{eq:tpc45-balanced-band}
\end{equation}
For this choice the endpoint exponents are
\begin{equation}
 \frac12=\frac{105000}{210000},
 \qquad
 1-\lambda_L=\frac{110021}{210000};
 \label{eq:tpc45-band-endpoint-exponents}
\end{equation}
their gap is
\begin{equation}
 \frac{5021}{210000}>0.
 \label{eq:tpc45-band-width-exponent}
\end{equation}
Constants in the dyadic endpoints cannot close this fixed power gap.

\begin{lemma}[Unique large prime and controlled cofactor]
\label{lem:tpc45-unique-balanced-cofactor}
Suppose that \(q\in\cP_R\) divides a nonzero target
\(n_{\alpha,j}\), and write
\begin{equation}
 n_{\alpha,j}=qr.
 \label{eq:tpc45-target-q-r}
\end{equation}
Then \(q\) is the only prime from \(\cP_R\) dividing that
target, and
\begin{equation}
 R\le r<\sqrt{N_+}.
 \label{eq:tpc45-balanced-r-range}
\end{equation}
Moreover
\begin{equation}
 s:=d_\alpha r
 \label{eq:tpc45-residual-label}
\end{equation}
is squarefree, \((q,s)=1\), and
\begin{equation}
 u_{\alpha,j}=qs.
 \label{eq:tpc45-u-q-s}
\end{equation}
\end{lemma}

\begin{proof}
Two primes larger than \(\sqrt{N_+}\) have product larger than every
target in the packet, which proves uniqueness.  The definition of the
upper endpoint of \(\cP_R\) gives
\[
 r=\frac{n_{\alpha,j}}q
 \ge \frac{N_-}{q}\ge R,
\]
whereas \(q>\sqrt{N_+}\) and \(n_{\alpha,j}\le N_+\) give the strict
upper bound in \eqref{eq:tpc45-balanced-r-range}.  Squarefreeness and
\eqref{eq:tpc45-source-target-coprime} show that \(q,r,d_\alpha\)
are pairwise coprime and squarefree.  The last two assertions follow.
\end{proof}

Since \(d_+=X^{10049/52500+o(1)}\) whereas
\(\sqrt{N_+}=X^{1/2+o(1)}\), one has
\begin{equation}
 d_+<\sqrt{N_+}
 \label{eq:tpc45-source-below-band}
\end{equation}
for all sufficiently large \(X\).  Hence no prime in \(\cP_R\) divides
a source factor \(d_\alpha\); band incidence in the fused label is
exactly the target incidence in \(n_{\alpha,j}\).

For the chosen \(R\), define each nonzero atom's descended label and
deterministic descent sign by
\begin{equation}
 (w_{\alpha,j},\eta_{\alpha,j})
 :=
 \begin{cases}
  (u_{\alpha,j}/q,-1),
   &q\in\cP_R,\ q\mid n_{\alpha,j},\\
  (u_{\alpha,j},1),&
   \text{if no prime of \(\cP_R\) divides \(n_{\alpha,j}\)}.
 \end{cases}
 \label{eq:tpc45-descended-label-sign}
\end{equation}
The first case is unambiguous by
\cref{lem:tpc45-unique-balanced-cofactor}, and there
\(w_{\alpha,j}=s=d_\alpha r\).

\begin{proposition}[Exact global-band character descent]
\label{prop:tpc45-character-descent}
Let \(\varepsilon\) denote the signs outside \(\cP_R\).  For
every nonzero equality atom,
\begin{equation}
 \boxed{
 \chi_{u_{\alpha,j}}
  (-\one_{\cP_R},\varepsilon)
 =\eta_{\alpha,j}\chi_{w_{\alpha,j}}(\varepsilon).}
 \label{eq:tpc45-exact-character-descent}
\end{equation}
At the all--minus point of the remaining signs,
\begin{equation}
 \eta_{\alpha,j}\mu(w_{\alpha,j})
 =\mu(u_{\alpha,j}).
 \label{eq:tpc45-exact-physical-sign}
\end{equation}
Thus the descended packet has exactly the original physical M\"obius
corner.
\end{proposition}

\begin{proof}
Untouched atoms give both identities trivially.  On a touched atom,
\(u=qs\), so
\[
 \chi_u(-\one_{\cP_R},\varepsilon)
 =-\chi_s(\varepsilon).
\]
This is \eqref{eq:tpc45-exact-character-descent}.  Since \(q\) and
\(s\) are coprime and squarefree,
\[
 -\mu(s)=\mu(q)\mu(s)=\mu(qs)=\mu(u),
\]
which proves \eqref{eq:tpc45-exact-physical-sign}.
\end{proof}

The use of the full-support quantities \(N_-\) and \(R\) in
\eqref{eq:tpc45-global-cofactor-band} is essential.  If a prime were selected
because one of its occurrences happened to have \(r\ge L_0\), another
occurrence of the same global prime could have a smaller cofactor.  The
uniform inequality \(r\ge R\) for every occurrence is what makes the
compression in the next section valid for the whole packet.  The clean
choice \(R=\ell_+\) is the first scale at which the fiber cost below is
subpolynomial without additional cancellation.
